% ===========================================================================
\section{Arcs, crosscuts and transfer}\label{sec:plane}
% ===========================================================================

\subsection{Arcs and curvilinear limits}

\begin{definition}[arcs]\label{def:arcs}
\lean{TJK.IsArc, TJK.IsArcAt, TJK.IsCrosscut}
A \emph{path} in $\R^n$ is a continuous $g:[0,1]\to\R^n$; a \emph{simple arc} is the image of an
injective path, and its \emph{endpoints} are $g(0),g(1)$. Let $\Omega\subseteq\R^n$ be open and
$p\in\partial\Omega$. An \emph{arc at $p$} (in $\Omega$) is a simple arc $\gamma$ with one endpoint $p$
and $\gamma\setminus\{p\}\subseteq\Omega$. A \emph{crosscut} of $\D$ is a simple arc whose endpoints
are distinct points of $C$ and whose other points lie in $\D$; a crosscut of $\Hp$ is defined in the
same way with $\R$ in place of $C$.
\end{definition}

\begin{definition}[curvilinear convergence, boundary functions, cluster sets]\label{def:ccset}
\lean{TJK.curvConvSet, TJK.IsBoundaryFunction, TJK.clusterSet}
\uses{def:arcs}
Let $\Omega$ be open, $B\subseteq\partial\Omega$, $Y$ a metric space and $f:\Omega\to Y$. For an arc
$\gamma$ at $p$, $f(z)\to q$ \emph{along $\gamma$} means $f(z)\to q$ as $z\to p$ with
$z\in\gamma\setminus\{p\}$. The \emph{set of curvilinear convergence} of $f$ (relative to $B$) is
\[
\CC(f)=\{p\in B:\ \text{there are an arc $\gamma$ at $p$ and $q\in Y$ with $f\to q$ along $\gamma$}\}.
\]
If $E\subseteq B$, a function $\varphi:E\to Y$ is a \emph{boundary function} for $f$ if for every
$p\in E$ there is an arc $\gamma$ at $p$ with $f\to\varphi(p)$ along $\gamma$. The \emph{cluster set}
of $f$ along $\gamma$ is $C(f,\gamma)=\{q:\exists z_n\in\gamma\setminus\{p\},\ z_n\to p,\ f(z_n)\to q\}$.
The pairs $(\Omega,B)$ used are $(\D,C)$, $(\Hp,\R)$ and $(Q^\circ,X)$, and in Section~\ref{sec:cube}
the open cube with its whole boundary.
\end{definition}

\begin{fact}[Tietze {\cite{Tietze}}]\label{fact:tietze-arc}
\lean{TJK.exists_arc_in_path}
\uses{def:arcs}
\source{Tietze 1919; also: every path-connected Hausdorff space is arc-connected}
If $g$ is a path in $\R^n$ with $g(0)\ne g(1)$, then $g([0,1])$ contains a simple arc with endpoints
$g(0)$ and $g(1)$.
\end{fact}

\begin{corollary}[paths may replace arcs]\label{cor:path-arc}
\lean{TJK.exists_arcAt_of_path, TJK.tendsto_of_nhdsWithin_le_map, TJK.clusterSet_subset_of_nhdsWithin_le_map}
\uses{fact:tietze-arc, def:ccset}
Let $\Omega$ be open, $p\in\partial\Omega$, and $g:[0,1]\to\R^n$ continuous with $g(1)=p$ and
$g([0,1))\subseteq\Omega$. Then $g([0,1])$ contains an arc $\gamma$ at $p$, and for every
$f:\Omega\to Y$: if $f(g(t))\to q$ as $t\to1$ then $f\to q$ along $\gamma$; and $C(f,\gamma)$ is
contained in the set of limits of $f(g(t_n))$ with $t_n\to1$.
\end{corollary}
\begin{proof}
Fact~\ref{fact:tietze-arc} gives a simple arc $\gamma\subseteq g([0,1])$ from $g(0)$ to $p$; since
$g(t)\ne p$ for $t<1$, $\gamma\setminus\{p\}\subseteq\Omega$. If $z_k\in\gamma\setminus\{p\}$ and
$z_k\to p$, write $z_k=g(t_k)$ with $t_k<1$; a subsequence of $t_k$ bounded away from $1$ would have a
limit $t^*<1$ with $g(t^*)=p$, which is impossible. So $t_k\to1$.
\end{proof}

\begin{lemma}[subarcs and exits]\label{lem:subarc}
\lean{TJK.IsArcAt.exists_subarc_subset, TJK.IsArcAt.tendsto_subarc, TJK.IsArcAt.clusterSet_subarc,
  TJK.IsArc.nhdsWithin_subarc, TJK.exists_exit_subpath, TJK.IsArcAt.exists_subarc_firstHit,
  TJK.IsArcAt.exists_subarc_strip}
\uses{def:ccset}
(a) If $\gamma$ is an arc at $p$ and $W$ is a neighbourhood of $p$, there is an arc $\gamma'\subseteq\gamma\cap
W$ at $p$; limits and cluster sets along $\gamma'$ and $\gamma$ coincide. If $f\to q$ along $\gamma$ and
$V\ni q$ is open, some such $\gamma'$ has $f(\gamma'\setminus\{p\})\subseteq V$.
(b) If $g$ is a path, $w=g(t_0)$, and $\abs{g(t_1)-w}\ge r>0$, there is a closed interval $J$ with
endpoints $t_0$ and some $t'$ between $t_0$ and $t_1$ such that $g(J)\subseteq\overline{S(r,w)}$ and
$\abs{g(t')-w}=r$; in particular $\diam g(J)\ge r$.
(c) If $\gamma$ is an arc at $p$ in $\Hp$ and $n$ is large enough that $\gamma$ meets $X_n$, then the subarc
of $\gamma$ from $p$ to its first point on $X_n$ is an arc at $p$ contained in $\overline{\Hp}_n$ with all
points other than its endpoints in $\Hp_n$.
\end{lemma}
\begin{proof}
Parametrize and take initial segments; for (b) let $t'$ be the first parameter between $t_0$ and $t_1$
with $\abs{g(t')-w}\ge r$, and use continuity.
\end{proof}

\begin{lemma}[cluster sets]\label{lem:cluster}
\lean{TJK.clusterSet_eq_iInter_subarc, TJK.clusterSet_nonempty, TJK.isCompact_clusterSet,
  TJK.isConnected_clusterSet, TJK.tendsto_iff_clusterSet_eq_singleton, TJK.clusterSet_comp}
\uses{def:ccset, lem:subarc}
Let $K$ be a compact metric space, $f:\Omega\to K$ continuous and $\gamma$ an arc at $p$. Then
$C(f,\gamma)=\bigcap_{\gamma'}\overline{f(\gamma'\setminus\{p\})}$, the intersection over the subarcs
$\gamma'$ of $\gamma$ at $p$; it is nonempty, compact and connected; and $f\to q$ along $\gamma$ iff
$C(f,\gamma)=\{q\}$. If $\Theta:K\to K'$ is continuous into a metric space, then
$C(\Theta\circ f,\gamma)=\Theta(C(f,\gamma))$.
\end{lemma}
\begin{proof}
The sets $\overline{f(\gamma'\setminus\{p\})}$ are compact, connected (images of half-open arcs) and
decrease as $\gamma'$ shrinks; a decreasing intersection of nonempty compact connected sets is
nonempty, compact and connected. The claim about limits is the usual compactness argument. For the last
claim: $\supseteq$ by continuity of $\Theta$; if $\Theta(f(z_n))\to q'$ with $z_n\to p$ along $\gamma$, a subsequence
of $f(z_n)$ converges in the compact $K$ to some $q\in C(f,\gamma)$, and $\Theta(q)=q'$.
\end{proof}

\begin{fact}[Tietze extension]\label{fact:tietze-ext}
\lean{TJK.tietze_bounded_complex}
\source{Mathlib: ContinuousMap.exists_extension / BoundedContinuousFunction extensions (apply to real and imaginary parts)}
A bounded continuous function from a closed subset of a metric space to $\C$ extends to a bounded continuous function on the
whole space.
\end{fact}

\begin{fact}[hyperspace]\label{fact:hyperspace}
\lean{TJK.compactSpace_nonemptyCompacts}
\source{Hausdorff, Mengenlehre p.~150; Mathlib: EMetric.NonemptyCompacts.compactSpace}
If $K$ is a compact metric space, the space $\mathcal C(K)$ of nonempty closed subsets of $K$ with the Hausdorff metric is
compact.
\end{fact}

\begin{fact}[Baire category theorem]\label{fact:baire-category}
\lean{TJK.baire_category}
\source{Mathlib: BaireSpace instances, dense_iInter_of_isOpen}
A complete metric space is not of the first category in itself; in it, a countable intersection of dense open sets is dense.
\end{fact}

\subsection{Crosscuts}

\begin{fact}[crosscut theorem]\label{fact:crosscut}
\lean{TJK.crosscut_theorem, TJK.closure_inter_circle_of_frontier, TJK.IsCrosscut.exists_angles}
\uses{def:arcs}
\source{Jordan curve theorem; Newman, Elements, Theorem 11.8 (the $\theta$-curve theorem)}
Let $\sigma$ be a crosscut of $\D$ with endpoints $a\ne b$, and let $A_1,A_2$ be the two closed arcs of
$C$ with endpoints $a,b$. Then $\D\setminus\sigma$ has exactly two components $U_1,U_2$, and the frontier
of $U_i$ is $\sigma\cup A_i$ ($i=1,2$). In particular $\overline{U_i}\cap C=A_i$ and
$\sigma\subseteq\overline{U_1}\cap\overline{U_2}$.
\end{fact}

\begin{lemma}[disjoint components have disjoint boundary arcs]\label{lem:disjoint-arcs}
\lean{TJK.disjoint_boundaryArcs}
\uses{fact:crosscut}
Let $\sigma,\tau$ be crosscuts of $\D$, $U$ a component of $\D\setminus\sigma$ and $V$ a component of
$\D\setminus\tau$, with $U\cap V=\emptyset$. Then the arcs $\overline U\cap C$ and $\overline V\cap C$ have
disjoint relative interiors.
\end{lemma}
\begin{proof}
Let $p$ be interior to both arcs. Then $p$ is not an endpoint of $\sigma$ or $\tau$, so $p\notin\sigma\cup\tau$,
and there is $r>0$ with $S(r,p)\cap(\sigma\cup\tau)=\emptyset$. The convex set $S(r,p)\cap\D$ lies in one
component of $\D\setminus\sigma$; as $p\in\overline U$ it meets $U$, so it lies in $U$. Likewise it lies in
$V$, contradicting $U\cap V=\emptyset$.
\end{proof}

\begin{definition}[charts]\label{def:charts}
\lean{TJK.cayley, TJK.chartH, TJK.chartD, TJK.stripExp, TJK.openSquare, TJK.baseX, TJK.gOpen,
  TJK.gClosed}
\begin{enumerate}[label=(\alph*)]
\item The Cayley map $\kappa(z)=(z-i)/(z+i)$.
\item $\Psi_\Hp\langle x,y\rangle=\tan(\pi x/2)+i\,y/(1-y)$ for $\abs x<1$, $0\le y<1$.
\item $\Psi_\D\langle x,y\rangle=1-1/w$ with $w=\tfrac12(1+y)+i\tan(\pi x/2)$, for $\abs x<1$, $0\le y<1$.
\item $\mathcal E_h(z)=e^{\pi z/h}$ for $h>0$.
\end{enumerate}
\end{definition}

\begin{lemma}[charts]\label{lem:charts}
\lean{TJK.cayley_isChart, TJK.tendsto_cayley_cobounded, TJK.tendsto_norm_cayleyInv,
  TJK.chartH_isChart, TJK.chartH_ofReal, TJK.strictMonoOn_tan_pi_mul_div_two, TJK.chartD_isChart,
  TJK.isClosed_gClosed, TJK.exists_ball_inter_disc_subset_gClosed, TJK.ofReal_not_mem_gClosed,
  TJK.stripExp_isChart, TJK.stripExp_ofReal, TJK.stripExp_add_I}
\uses{def:charts}
(a) $\kappa$ is a homeomorphism of $\Hp\cup\R$ onto $(\D\cup C)\setminus\{1\}$ mapping $\Hp$ onto $\D$
and $\R$ onto $C\setminus\{1\}$; $\abs{\kappa(z)-1}\to0$ iff $\abs z\to\infty$.
(b) $\Psi_\Hp$ is a homeomorphism of $Q^\circ\cup X$ onto $\Hp\cup\R$ mapping $X$ onto $\R$
increasingly.
(c) $\Psi_\D$ is a homeomorphism of $Q^\circ\cup X$ onto $G^\circ\cup(C\setminus\{1\})$, where
$G^\circ=\{z\in\D:\operatorname{Re}(1-z)^{-1}<1\}$, mapping $X$ onto $C\setminus\{1\}$. The set
$G''=\{z\in\D:\operatorname{Re}(1-z)^{-1}\le3/4\}$ is closed in $\D$, contains a neighbourhood in $\D$
of each point of $C\setminus\{1\}$, and is disjoint from the radius $[0,1)$.
(d) $\mathcal E_h$ maps $\{0\le\operatorname{Im}z\le h\}$ homeomorphically onto
$(\Hp\cup\R)\setminus\{0\}$, the line $\operatorname{Im}z=0$ increasingly onto $(0,\infty)$ and the line
$\operatorname{Im}z=h$ decreasingly onto $(-\infty,0)$.
\end{lemma}
\begin{proof}
(a) $\abs{\kappa(z)}<1$ iff $\abs{z-i}<\abs{z+i}$ iff $\operatorname{Im}z>0$; the inverse is
$w\mapsto i(1+w)/(1-w)$. (b) is a product of homeomorphisms $(-1,1)\to\R$ and $[0,1)\to[0,\infty)$.
(c) $z\mapsto(1-z)^{-1}$ maps $\D$ onto $\{\operatorname{Re}w>\frac12\}$ and $C\setminus\{1\}$ onto
$\{\operatorname{Re}w=\frac12\}$ (since $\abs{1-1/w}<1\iff\abs{w-1}<\abs w$); $\Psi_\D$ is the inverse
of this map composed with a homeomorphism of $Q^\circ\cup X$ onto $\{\frac12\le\operatorname{Re}w<1\}$.
Near a point $x\in C\setminus\{1\}$, $\operatorname{Re}(1-z)^{-1}$ is close to $\frac12$; on $[0,1)$ it is
$\ge1$. (d) is elementary.
\end{proof}

\begin{lemma}[crosscuts of the half-plane]\label{lem:crosscut-H}
\lean{TJK.crosscutH_theorem}
\uses{fact:crosscut, lem:charts}
Let $\sigma$ be a crosscut of $\Hp$ with endpoints $a<b$. Then $\Hp\setminus\sigma$ has exactly two
components, a bounded one $V$ with $\overline V\cap\R=[a,b]$, and an unbounded one $U$ with
$\overline U\cap\R=\R\setminus(a,b)$; moreover $\sigma\subseteq\overline U\cap\overline V$ and
$U\supseteq\{\operatorname{Im}z>\max_\sigma\operatorname{Im}\}$.
\end{lemma}
\begin{proof}
$\kappa(\sigma)$ is a crosscut of $\D$ with endpoints in $C\setminus\{1\}$
(Lemma~\ref{lem:charts}(a)); apply Fact~\ref{fact:crosscut}. The component whose closed arc avoids $1$
has closure a compact subset of $\overline\D\setminus\{1\}$, so its preimage $V$ is bounded, and $\kappa$
maps $[a,b]$ onto that arc. The half-plane above $\sigma$ is connected, unbounded and disjoint from
$\sigma$, so it lies in $U$.
\end{proof}

\begin{lemma}[order of disjoint arcs in a strip]\label{lem:strip}
\lean{TJK.strip_order}
\uses{lem:crosscut-H, lem:charts}
Let $h>0$, $S=\{0<\operatorname{Im}z<h\}$, and let $\alpha,\alpha'$ be disjoint simple arcs, $\alpha$ from
$x\in\R$ to $p$ with $\operatorname{Im}p=h$, $\alpha'$ from $x'\in\R$ to $p'$ with
$\operatorname{Im}p'=h$, all other points in $S$. If $x<x'$ then $\operatorname{Re}p<\operatorname{Re}p'$.
\end{lemma}
\begin{proof}
By Lemma~\ref{lem:charts}(d), $\mathcal E_h(\alpha)$ is a crosscut of $\Hp$ from $u=e^{\pi x/h}$ to
$v=-e^{\pi\operatorname{Re}p/h}$, and $\mathcal E_h(\alpha')$ one from $u'>u$ to
$v'=-e^{\pi\operatorname{Re}p'/h}\ne v$. By Lemma~\ref{lem:crosscut-H} the bounded component $V$ of
$\Hp\setminus\mathcal E_h(\alpha)$ has $\overline V\cap\R=[v,u]$. The connected set
$\mathcal E_h(\alpha')\setminus\{u',v'\}$ lies in one component; its closure contains $u'\notin[v,u]$, so
it lies in the unbounded component $U$, and then $v'\in\overline U\cap\R=\R\setminus(v,u)$. As $v'<0<u$
and $v'\ne v$, $v'<v$, i.e.\ $\operatorname{Re}p'>\operatorname{Re}p$.
\end{proof}

\begin{lemma}[geometry near a boundary point]\label{lem:disc-geometry}
\lean{TJK.convex_disc_inter_ball, TJK.isConnected_disc_diff_closedBall, TJK.shortArc_subset_ball,
  TJK.exists_shortArc, TJK.norm_expI_sub_le, TJK.sum_le_two_pi_of_pairwiseDisjoint, TJK.card_mul_le_two_pi}
Let $x\in C$ and $0<\delta<2$.
(a) $\D\cap S(\delta,x)$ is convex.
(b) $\D\setminus\overline{S(\delta,x)}$ is connected and has $-x$ in its closure.
(c) If $x'\in C$ and $\abs{x'-x}<\delta$, the shorter closed arc of $C$ from $x$ to $x'$ lies in $S(\delta,x)$.
Any arc of $C$ between $w_1,w_2$ has length at least $\abs{w_1-w_2}$.
(d) If $A_1,\dots,A_m$ are closed arcs of $C$ with pairwise disjoint relative interiors, their lengths add up to at most
$2\pi$. (The length of an arc is the Lebesgue measure of the corresponding set of angles in $[0,2\pi)$.) In particular,
if the endpoints of each $A_i$ are at distance $\ge\eps$, then $m\eps\le2\pi$.
In Lean a closed arc is $\{e^{it}:\theta\le t\le\theta+\ell\}$ with $0\le\ell\le2\pi$ (\texttt{circleArc}), its interior
is the image of $(\theta,\theta+\ell)$, and its length is $\ell$; (d) is proved by pushing the open angle intervals to
$\R/2\pi\Z$ with its Haar measure.
\end{lemma}
\begin{proof}
(a) Intersection of convex sets. (b) For $z\in\D$ the convex function $t\mapsto\abs{(1-t)z-tx-x}^2$ has
derivative $2(1-\abs z^2)>0$ at $t=0$, hence increases on $[0,1]$; so the segment $[z,-x)$ stays in
$\D\setminus\overline{S(\delta,x)}$ and ends in the convex set $\D\cap S(2-\delta,-x)$, which lies in
$\D\setminus\overline{S(\delta,x)}$. (c) On an arc of angle at most $\pi$ the chord from an endpoint
increases with the angle; chords are shorter than arcs. (d) The interiors are disjoint measurable sets of
angles and the endpoints are null.
\end{proof}

\begin{lemma}\label{lem:open-components}
\lean{TJK.countable_connectedComponents_of_isOpen}
An open subset of $\R^n$ has countably many components; they are open and path-connected, and a
connected subset of the open set lies in one of them.
\end{lemma}
\begin{proof}
$\R^n$ is locally path-connected and second countable.
\end{proof}

\subsection{Transfer}

\begin{lemma}[transfer]\label{lem:transfer}
\lean{TJK.IsChart, TJK.transfer_isArcAt, TJK.transfer_tendsto, TJK.transfer_curvConvSet,
  TJK.transfer_isBoundaryFunction, TJK.transfer_clusterSet}
\uses{def:ccset}
Let $\Omega,\Omega'$ be open, $B\subseteq\partial\Omega$, $B'\subseteq\partial\Omega'$, and let
$\Psi:\Omega\cup B\to\Omega'\cup B'$ be a homeomorphism with $\Psi(\Omega)=\Omega'$. For $f:\Omega'\to Y$:
$\gamma$ is an arc at $p\in B$ iff $\Psi(\gamma)$ is an arc at $\Psi(p)$; $f\circ\Psi\to q$ along
$\gamma$ iff $f\to q$ along $\Psi(\gamma)$, and the cluster sets agree; hence
$\CC(f\circ\Psi)=\Psi^{-1}(\CC(f))$, and $\varphi$ is a boundary function for $f$ on $E'$ iff
$\varphi\circ\Psi$ is one for $f\circ\Psi$ on $\Psi^{-1}(E')$.
\end{lemma}
\begin{proof}
Arcs at $p$ lie in $\Omega\cup\{p\}$, where $\Psi$ is a homeomorphism.
\end{proof}

\begin{lemma}[locality]\label{lem:locality}
\lean{TJK.curvConv_local}
\uses{def:ccset, lem:subarc}
Let $p\in\partial\Omega\cap\partial\Omega'$ and suppose $W\cap\Omega=W\cap\Omega'$ for an open $W\ni p$. If
$f=f'$ on $W\cap\Omega$, then $f$ has a limit $q$ along some arc at $p$ in $\Omega$ iff $f'$ has the limit
$q$ along some arc at $p$ in $\Omega'$.
\end{lemma}
\begin{proof}
Lemma~\ref{lem:subarc}(a).
\end{proof}
